\section{关键质量场景}

非功能需求给出了系统级目标，表~\ref{tab:quality-scenarios} 进一步将其落实为可观察的运行场景。表中的“设计目标”表示尚未形成稳定实测记录，不能作为当前版本已经达到的性能结论；“已有自动化验证”表示仓库中存在对应测试证据。

\begingroup
\footnotesize
\setlength{\tabcolsep}{3pt}
\begin{longtable}{@{}p{1.1cm} p{1.8cm} p{2.1cm} p{3.2cm} p{1.8cm} p{1.7cm}@{}}
\caption{关键质量场景}
\label{tab:quality-scenarios}\\
\toprule
\textbf{属性} & \textbf{场景} & \textbf{刺激与环境} & \textbf{系统响应} & \textbf{度量} & \textbf{验证状态} \\
\midrule
\endfirsthead
\multicolumn{6}{c}{\small 表~\thetable\quad 关键质量场景（续）} \\
\toprule
\textbf{属性} & \textbf{场景} & \textbf{刺激与环境} & \textbf{系统响应} & \textbf{度量} & \textbf{验证状态} \\
\midrule
\endhead
\midrule
\multicolumn{6}{r}{\small 接下页} \\
\endfoot
\bottomrule
\endlastfoot
性能 & 千笔账单导入 & 登录用户上传不超过 1000 笔交易的微信或支付宝账单 & 完成平台识别、解析、规范化、去重并持续更新批次进度；模型等待不应阻塞进度查询 & 解析阶段目标 10 秒内；进度可查询 & 设计目标；流程有自动化测试 \\
\hline
可靠性 & 外部模型超时 & 后台任务处理器调用外部模型超过配置时间或请求失败 & 保留交易记录，增加重试计数并按指数退避重新调度；达到上限后记录失败原因 & 导入主流程不因单笔模型超时整体失败 & 已有自动化验证 \\
\hline
安全性 & 跨用户资源访问 & 已登录用户修改交易、批次或报表编号，尝试访问他人资源 & 根据当前用户与资源归属拒绝请求，不返回目标资源内容 & 返回 403 或 404；无数据泄露 & 已有接口自动化验证 \\
\hline
可靠性 & 部分交易处理失败 & 单个文件中部分记录解析或分类失败 & 成功记录继续保存，错误信息保留，批次显示为部分失败 & 成功数据可继续使用；失败原因可追踪 & 已有自动化验证 \\
\hline
可扩展性 & 新增账单平台 & 开发者增加一种新的平台账单格式 & 在平台识别与解析模块中增加相应解析规则，并输出统一交易结构，无需修改后续分类和分析接口 & 修改集中在解析模块及其测试 & 由现有解析机制验证 \\
\hline
可维护性 & 切换分类服务 & 用户或部署人员切换规则、外部接口或本地模型配置 & 通过统一分类结果保持导入、校正和分析模块接口不变 & 不修改调用方接口；失败可降级或重试 & 已有分类功能测试 \\
\hline
性能 & 大数据量财务看板 & 用户选择较长日期范围或家庭组织聚合口径 & 后端执行按用户范围和日期过滤的聚合查询，并返回摘要、分布和排行 & 普通交互目标 P95 小于 1 秒 & 设计目标，尚缺专项压测 \\
\end{longtable}
\endgroup
